\documentclass[12pt]{article}
\usepackage[utf8]{inputenc}
\usepackage[T1]{fontenc}
\usepackage{hyperref}
\title{Talise Language Resources for K-12 Curriculum Development}
\author{Compiled by Claude Sonnet 5 (claude-sonnet-5) for Andr\'as Kornai}
\date{}
\begin{document}
\maketitle

\section{Introduction}
Talise (ISO 639-3 tlr; Glottocode tali1259) is a Southeast Solomonic (Oceanic, Austronesian) language of Guadalcanal, Solomon Islands. It is also called Tolo or Talisi. Search results give speaker numbers far below the figure in the task description: Wikipedia cites about 13,000 (1999), and ASJP lists 12,500 for Moli Talise and 19,300 for Koo Talise. Some sources treat Talise as one language and others as an umbrella for Poleo, Koo, Malagheti, Moli and Tolo. Little material exists, and I found no school materials.

\section{Primary Educational Resources}
I found no primers, textbooks or school programmes for Talise in search results.

\section{Linguistic and Reference Materials}
\begin{itemize}
\item Wikipedia, Talise language (overview and phonology): \url{https://en.wikipedia.org/wiki/Talise_language}
\item Omniglot, Talise writing and alphabet: \url{https://www.omniglot.com/writing/talise.htm}
\item Rosetta Project page for Talise (lists grammar and phonology documents; I did not open them): \url{https://rosettapanglossia.longnow.org/index.php?title=Tlr}
\item Glottolog reference: Crowley, Susan Smith 1986 (the Tolo Dictionary, reported by search as the main dictionary reference; I did not see the text itself): \url{https://glottolog.org/resource/reference/id/27644}
\item Glottolog reference: Capell, Arthur 1930, a grammar sketch of the language of Inakona, Guadalcanal: \url{https://glottolog.org/resource/reference/id/115284}
\end{itemize}

\section{Religious and Community Resources}
None found in search results.

\section{Digital and Audio Resources}
\begin{itemize}
\item ASJP wordlist, Moli Talise: \url{https://asjp.clld.org/languages/MOLI_TALISE}
\item ASJP wordlist, Koo Talise: \url{https://asjp.clld.org/languages/KOO_TALISE}
\end{itemize}

\section{Curriculum Development Context}
Documentation is limited to a 1986 dictionary, an older grammar sketch and wordlists. Any curriculum would need new materials built with the community. Dialect variation and the uncertain speaker numbers should be settled first.

\end{document}